\chapter{Twierdzenia o wartości średniej}
\label{ma1:ch:mean-value-theorems}

Pochodna opisuje zachowanie funkcji w pojedynczym punkcie. Twierdzenia o wartości średniej pozwalają przejść od tej informacji lokalnej do informacji o przyroście funkcji na całym przedziale. Podstawą dowodów jest istnienie ekstremów funkcji ciągłej na przedziale domkniętym i ograniczonym.

\section{Twierdzenie Fermata}
\begin{definicja}
\label{ma1:def:local-extremum}
Niech $f:I\to\R$, gdzie $I$ jest przedziałem, i niech $c\in I$. Mówimy, że $f$ ma w $c$ maksimum lokalne, jeżeli istnieje $\delta>0$ takie, że $f(x)\le f(c)$ dla wszystkich $x\in I\cap(c-\delta,c+\delta)$. Minimum lokalne definiujemy przez przeciwną nierówność. Ekstremum lokalne jest maksimum lub minimum lokalnym. Jeżeli nierówność jest ostra dla wszystkich $x\ne c$ w tym otoczeniu, ekstremum nazywamy ścisłym.
\end{definicja}
\begin{twierdzenie}[Fermata]
\label{ma1:thm:fermat}
Jeśli $c\in\operatorname{int}I$, funkcja $f:I\to\R$ ma w $c$ ekstremum lokalne i jest różniczkowalna w $c$, to $f'(c)=0$.
\end{twierdzenie}
\begin{proof}
Załóżmy najpierw, że $c$ jest punktem maksimum lokalnego. Dla dostatecznie małych $h\ne0$ mamy $c+h\in I$ i $f(c+h)-f(c)\le0$. Zatem
\[
 \frac{f(c+h)-f(c)}h\le0 \quad(h>0),\qquad
 \frac{f(c+h)-f(c)}h\ge0 \quad(h<0).
\]
Obie granice jednostronne istnieją i są równe $f'(c)$, więc $f'(c)\le0$ i $f'(c)\ge0$. Stąd $f'(c)=0$. Jeśli $f$ ma w $c$ minimum lokalne, funkcja $-f$ ma tam maksimum lokalne i $(-f)'(c)=-f'(c)$, co daje tę samą tezę.
\end{proof}
Punkt wewnętrzny, w którym istnieje pochodna równa zeru, nazywamy punktem stacjonarnym. Twierdzenie Fermata daje warunek konieczny, lecz nie wystarczający istnienia ekstremum.
\begin{przyklad}
\label{ma1:ex:fermat-assumptions}
Funkcja $x^3$ ma pochodną równą zeru w $0$, ale dla $h>0$ mamy $h^3>0$, a dla $h<0$ mamy $h^3<0$, więc w $0$ nie ma ekstremum lokalnego. Funkcja $|x|$ ma minimum w $0$, ale nie jest tam różniczkowalna. Funkcja $x\mapsto x$ na $[0,1]$ ma minimum w $0$ i maksimum w $1$, a jej odpowiednie pochodne jednostronne wynoszą $1$. Założenie o punkcie wewnętrznym jest zatem istotne.
\end{przyklad}
W wersji jednostronnej, przy założeniu istnienia odpowiedniej skończonej pochodnej, w maksimum na lewym końcu przedziału mamy $f'_+(a)\le0$, a w maksimum na prawym końcu $f'_-(b)\ge0$. Dla minimum znaki są przeciwne. Wynika to bezpośrednio ze znaków ilorazu różnicowego w dowodzie.

\section{Twierdzenie Rolle'a}
\begin{twierdzenie}[Rolle'a]
\label{ma1:thm:rolle}
Niech $a<b$. Jeżeli $f:[a,b]\to\R$ jest ciągła na $[a,b]$, różniczkowalna na $(a,b)$ i $f(a)=f(b)$, to istnieje $c\in(a,b)$ takie, że $f'(c)=0$.
\end{twierdzenie}
\begin{proof}
Z twierdzenia \ref{ma1:cor:interval-extrema} funkcja osiąga minimum $m$ i maksimum $M$ na $[a,b]$. Jeśli $m=M$, funkcja jest stała, więc $f'(c)=0$ dla każdego $c\in(a,b)$. Niech $m<M$ i $v=f(a)=f(b)$. Ponieważ $m\le v\le M$, co najmniej jedna z nierówności $m<v$ lub $v<M$ jest ostra. Jeśli $m<v$, punkt osiągnięcia minimum nie może być żadnym z końców, więc istnieje $c\in(a,b)$, w którym $f(c)=m$. Jeśli $v<M$, analogicznie wybieramy wewnętrzny punkt osiągnięcia maksimum. W obu przypadkach stosujemy twierdzenie \ref{ma1:thm:fermat}.
\end{proof}
\begin{center}
\begin{tikzpicture}[x=3cm,y=2.1cm,>=Stealth]
 \draw[->] (-.12,0)--(2.25,0) node[right] {$x$};
 \draw[->] (0,-.12)--(0,1.75) node[above] {$y$};
 \draw[blue!65!black,thick,domain=.2:1.8,samples=100] plot (\x,{1.45-(\x-1)^2});
 \draw[green!40!black,dashed,thick] (.55,1.45)--(1.45,1.45);
 \draw[gray,dotted] (.2,0)--(.2,.81)--(1.8,.81)--(1.8,0);
 \draw[gray,dotted] (1,0)--(1,1.45);
 \fill (.2,.81) circle (1.4pt); \fill (1.8,.81) circle (1.4pt);
 \fill[green!40!black] (1,1.45) circle (1.4pt);
 \node[below] at (.2,0) {$a$};
 \node[below] at (1,0) {$c$};
 \node[below] at (1.8,0) {$b$};
 \node[left,font=\small] at (0,.81) {$f(a)$};
 \node[above,font=\small,green!40!black] at (1,1.45) {$f'(c)=0$};
\end{tikzpicture}

\small Rysunek: równość wartości na końcach przedziału wymusza istnienie stycznej poziomej w pewnym punkcie wewnętrznym.
\end{center}
\begin{przyklad}
\label{ma1:ex:rolle-assumptions}
Każde z założeń jest istotne dla ogólnej tezy.
\begin{itemize}
 \item Jeśli $f(0)=1$ i $f(x)=x$ dla $0<x\le1$, to $f(0)=f(1)$, funkcja jest różniczkowalna na $(0,1)$ i $f'(x)=1$, ale nie jest ciągła w $0$.
 \item Funkcja $f(x)=|x|$ na $[-1,1]$ jest ciągła i ma równe wartości na końcach, lecz jej pochodna w $0$ nie istnieje, a poza zerem przyjmuje wyłącznie wartości $-1$ i $1$.
 \item Funkcja $f(x)=x$ na $[0,1]$ jest ciągła i różniczkowalna, ale wartości na końcach są różne i pochodna nigdzie nie jest zerowa.
\end{itemize}
\end{przyklad}

\section{Twierdzenie Lagrange'a}
\begin{twierdzenie}[Lagrange'a o wartości średniej]
\label{ma1:thm:lagrange}
Jeżeli $a<b$, $f:[a,b]\to\R$ jest ciągła na $[a,b]$ i różniczkowalna na $(a,b)$, to istnieje $c\in(a,b)$ takie, że
\[
 f(b)-f(a)=f'(c)(b-a).
\]
\end{twierdzenie}
\begin{proof}
Niech $m=(f(b)-f(a))/(b-a)$ oraz $F(x)=f(x)-f(a)-m(x-a)$. Funkcja $F$ jest ciągła na $[a,b]$ i różniczkowalna na $(a,b)$ jako różnica funkcji o tych własnościach. Mamy $F(a)=0$ i $F(b)=f(b)-f(a)-m(b-a)=0$. Z twierdzenia \ref{ma1:thm:rolle} istnieje $c\in(a,b)$ takie, że $0=F'(c)=f'(c)-m$. Stąd $f'(c)=m$, co daje tezę po pomnożeniu przez $b-a$.
\end{proof}
Geometrycznie oznacza to istnienie stycznej równoległej do siecznej łączącej końce wykresu nad przedziałem. W szczególności pochodna przyjmuje wartość średniego ilorazu różnicowego. Twierdzenie nie orzeka, że taki punkt jest jedyny.
\begin{center}
\begin{tikzpicture}[x=3.2cm,y=1.7cm,>=Stealth]
 \draw[->] (-.12,0)--(2.2,0) node[right] {$x$};
 \draw[->] (0,-.1)--(0,3) node[above] {$y$};
 \draw[blue!65!black,thick,domain=.2:1.8,samples=100] plot (\x,{.7*\x*\x+.35});
 \draw[orange!85!black,thick] (.2,.378)--(1.8,2.618);
 \draw[green!40!black,dashed,thick] (.5,.35)--(1.65,1.96);
 \draw[gray,dotted] (.2,0)--(.2,.378);
 \draw[gray,dotted] (1,0)--(1,1.05);
 \draw[gray,dotted] (1.8,0)--(1.8,2.618);
 \fill (.2,.378) circle (1.4pt); \fill (1.8,2.618) circle (1.4pt);
 \fill[green!40!black] (1,1.05) circle (1.4pt);
 \node[below] at (.2,0) {$a$}; \node[below] at (1,0) {$c$};
 \node[below] at (1.8,0) {$b$};
 \node[orange!85!black,above left,font=\small] at (1.25,1.848) {sieczna};
 \node[green!40!black,below right,font=\small] at (1.5,1.75) {styczna};
\end{tikzpicture}

\small Rysunek: twierdzenie Lagrange'a dla $f(x)=0{,}7x^2+0{,}35$ na $[0{,}2,1{,}8]$. Sieczna i styczna w $c=1$ mają ten sam współczynnik kierunkowy $1{,}4$.
\end{center}
\begin{przyklad}
\label{ma1:ex:lagrange-witnesses}
Dla $f(x)=x^2$ na $[a,b]$ iloraz różnicowy wynosi $a+b$, więc równanie $2c=a+b$ wyznacza dokładnie jeden punkt $c=(a+b)/2\in(a,b)$. Dla $f(x)=x^3$ na $[-1,1]$ iloraz różnicowy wynosi $1$, więc równanie $3c^2=1$ ma dwa rozwiązania $c=\pm1/\sqrt3$, oba wewnątrz przedziału. Dla funkcji afinicznej każdy punkt wewnętrzny spełnia tezę.
\end{przyklad}

\section{Twierdzenie Cauchy'ego}
\begin{twierdzenie}[Cauchy'ego o wartości średniej]
\label{ma1:thm:cauchy-mvt}
Niech $a<b$ i niech $f,g:[a,b]\to\R$ będą ciągłe na $[a,b]$ oraz różniczkowalne na $(a,b)$. Istnieje $c\in(a,b)$ takie, że
\[
 (f(b)-f(a))g'(c)=(g(b)-g(a))f'(c).
\]
\end{twierdzenie}
\begin{proof}
Oznaczmy $A=f(b)-f(a)$ i $B=g(b)-g(a)$. Funkcja $H(x)=B(f(x)-f(a))-A(g(x)-g(a))$ jest ciągła na $[a,b]$ i różniczkowalna na $(a,b)$. Mamy $H(a)=0$ i $H(b)=BA-AB=0$. Twierdzenie Rolle'a daje $c\in(a,b)$ z $H'(c)=Bf'(c)-Ag'(c)=0$. Jest to żądana równość. Dowód nie wymaga dzielenia przez $A$, $B$ ani przez pochodne.
\end{proof}
\begin{wniosek}
\label{ma1:cor:cauchy-ratio}
Jeżeli dodatkowo $g'(x)\ne0$ dla każdego $x\in(a,b)$, to $g(b)\ne g(a)$ i istnieje $c\in(a,b)$ takie, że
\[
 \frac{f(b)-f(a)}{g(b)-g(a)}=\frac{f'(c)}{g'(c)}.
\]
\end{wniosek}
\begin{proof}
Gdyby $g(a)=g(b)$, twierdzenie Rolle'a zastosowane do $g$ dawałoby punkt $d\in(a,b)$ z $g'(d)=0$, wbrew założeniu. Oba mianowniki są zatem niezerowe dla punktu $c$ z twierdzenia \ref{ma1:thm:cauchy-mvt}, więc wolno podzielić jego równość przez ich iloczyn.
\end{proof}
Dla $g(x)=x$ twierdzenie Cauchy'ego daje twierdzenie Lagrange'a. Geometrycznie, w wersji ilorazowej (przy założeniach wniosku \ref{ma1:cor:cauchy-ratio}), wektor prędkości $(g'(c),f'(c))$ krzywej parametrycznej $x\mapsto(g(x),f(x))$ jest równoległy do wektora łączącego jej końce: $(g(b)-g(a),f(b)-f(a))$. Ogólna równość Cauchy'ego wyraża znikanie wyznacznika tych dwóch wektorów; może zachodzić także wtedy, gdy któryś jest zerowy.
\begin{przyklad}
\label{ma1:ex:cauchy-degenerate}
Dla $f(x)=x$, $g(x)=x^2$ na $[-1,1]$ mamy $A=2$, $B=0$, więc równość Cauchy'ego zachodzi w $c=0$, ale iloraz $A/B$ nie jest określony, podobnie jak $f'(c)/g'(c)$. Nie wolno zastępować twierdzenia bezwarunkowo wersją ilorazową. Dla $f(x)=x^2$, $g(x)=x^3$ na $[1,2]$ mamy $A=3$, $B=7$, a $f'(c)/g'(c)=2/(3c)$. Punkt $c=14/9\in(1,2)$ spełnia równość ilorazów.
\end{przyklad}

\section{Znak pochodnej i monotoniczność}
\begin{twierdzenie}
\label{ma1:thm:derivative-monotonicity}
Niech $I$ będzie przedziałem o niepustym wnętrzu, a $f:I\to\R$ niech będzie ciągła na $I$ i różniczkowalna w każdym punkcie $\operatorname{int}I$. Wtedy:
\begin{enumerate}
 \item $f$ jest niemalejąca na $I$ wtedy i tylko wtedy, gdy $f'(x)\ge0$ dla każdego $x\in\operatorname{int}I$;
 \item $f$ jest nierosnąca na $I$ wtedy i tylko wtedy, gdy $f'(x)\le0$ dla każdego $x\in\operatorname{int}I$;
 \item jeśli $f'(x)>0$ dla każdego $x\in\operatorname{int}I$, to $f$ jest ściśle rosnąca; jeśli $f'(x)<0$ dla każdego takiego $x$, to $f$ jest ściśle malejąca;
 \item $f$ jest stała na $I$ wtedy i tylko wtedy, gdy $f'(x)=0$ dla każdego $x\in\operatorname{int}I$.
\end{enumerate}
\end{twierdzenie}
\begin{proof}
Jeżeli $f$ jest niemalejąca i $x\in\operatorname{int}I$, to dla dostatecznie małych $h>0$ licznik $f(x+h)-f(x)$ jest nieujemny, a dla $h<0$ jest niedodatni. W obu przypadkach iloraz różnicowy jest nieujemny, więc jego granica jest nieujemna.

Odwrotnie, niech $f'\ge0$ na $\operatorname{int}I$. Dla dowolnych $u,v\in I$, $u<v$, cały przedział $[u,v]$ zawiera się w $I$, a $(u,v)\subseteq\operatorname{int}I$. Ograniczenie $f$ spełnia założenia twierdzenia Lagrange'a. Istnieje zatem $c\in(u,v)$ takie, że $f(v)-f(u)=f'(c)(v-u)\ge0$. Ponieważ $u,v$ są dowolne, $f$ jest niemalejąca. Przy założeniu $f'>0$ ta sama równość daje $f(v)>f(u)$, czyli ścisły wzrost. Stosując oba wyniki do $-f$, otrzymujemy odpowiednie twierdzenia o nierosnącym i ściśle malejącym charakterze. Jeśli $f'=0$, pierwsze dwie części mówią, że $f$ jest jednocześnie niemalejąca i nierosnąca, więc jest stała. Pochodna funkcji stałej jest oczywiście zerowa.
\end{proof}
Ścisły wzrost nie implikuje dodatniości pochodnej w każdym punkcie: funkcja $x^3$ jest ściśle rosnąca, lecz $f'(0)=0$. Istotnie, dla $u<v$ różnica $v^3-u^3=(v-u)(v^2+uv+u^2)$ jest dodatnia, ponieważ $v^2+uv+u^2=(v+u/2)^2+3u^2/4>0$ dla $u\ne v$.

\begin{wniosek}
\label{ma1:cor:strict-monotone-nonflat}
Przy założeniach twierdzenia \ref{ma1:thm:derivative-monotonicity} i warunku $f'\ge0$ funkcja $f$ jest ściśle rosnąca wtedy i tylko wtedy, gdy w każdym niepustym otwartym przedziale $(u,v)\subseteq I$ istnieje $c\in(u,v)$ z $f'(c)>0$.
\end{wniosek}
\begin{proof}
Jeżeli $f$ jest ściśle rosnąca, wybieramy $u<p<q<v$ i stosujemy twierdzenie Lagrange'a do $[p,q]$. Dostajemy $f'(c)=(f(q)-f(p))/(q-p)>0$ w pewnym $c\in(p,q)$. Odwrotnie, $f'\ge0$ daje niemalejący charakter. Gdyby dla pewnych $u<v$ w $I$ zachodziło $f(u)=f(v)$, to dla każdego $x\in[u,v]$ mielibyśmy $f(u)\le f(x)\le f(v)=f(u)$. Funkcja byłaby stała na $[u,v]$, a więc jej pochodna znikałaby na $(u,v)$, wbrew warunkowi.
\end{proof}
\begin{wniosek}
\label{ma1:cor:equal-derivatives}
Jeżeli $f,g$ są ciągłe na przedziale $I$ o niepustym wnętrzu, różniczkowalne w jego wnętrzu i $f'=g'$ w tym wnętrzu, to istnieje $C\in\R$ takie, że $f=g+C$ na $I$.
\end{wniosek}
\begin{proof}
Funkcja $f-g$ spełnia te same założenia regularności i ma pochodną równą zeru. Z twierdzenia \ref{ma1:thm:derivative-monotonicity} jest stała.
\end{proof}
Założenie o przedziale nie może być pominięte. Na $\R\setminus\{0\}$ funkcja równa $0$ dla $x<0$ i $1$ dla $x>0$ jest różniczkowalna i ma pochodną zerową, lecz nie jest stała na całej dziedzinie.

\section{Oszacowania przyrostów i zastosowania}
\begin{twierdzenie}
\label{ma1:thm:mvt-increment-bounds}
Niech $f$ będzie ciągła na przedziale $I$ o niepustym wnętrzu i różniczkowalna w jego wnętrzu. Jeżeli $m\le f'(x)\le M$ dla każdego $x\in\operatorname{int}I$, gdzie $m,M\in\R$, to dla wszystkich $u<v$ w $I$ zachodzi
\[
 m(v-u)\le f(v)-f(u)\le M(v-u).
\]
W szczególności, jeśli $|f'(x)|\le L$ dla każdego $x\in\operatorname{int}I$, gdzie $L\ge0$, to $f$ jest funkcją Lipschitza ze stałą $L$ na $I$.
\end{twierdzenie}
\begin{proof}
Dla $u<v$ w $I$ twierdzenie Lagrange'a daje $f(v)-f(u)=f'(c)(v-u)$ z $c\in(u,v)$. Mnożenie nierówności dla $f'(c)$ przez dodatnie $v-u$ zachowuje ich kierunek. Przy $m=-L$, $M=L$ otrzymujemy $|f(v)-f(u)|\le L(v-u)$. Dla $u>v$ zamieniamy role punktów, a dla $u=v$ obie strony są zerowe. To jest warunek Lipschitza.
\end{proof}
Odwrotnie, jeżeli $f$ jest funkcją Lipschitza ze stałą $L$ i jest różniczkowalna w punkcie wewnętrznym $x$, to $|(f(x+h)-f(x))/h|\le L$ dla każdego dostatecznie małego $h\ne0$. Przejście do granicy daje $|f'(x)|\le L$. W połączeniu z twierdzeniem oznacza to równoważność ograniczoności pochodnej i warunku Lipschitza dla funkcji ciągłych na przedziale i różniczkowalnych w jego wnętrzu. Sama własność Lipschitza nie daje jednak różniczkowalności we wszystkich punktach, co pokazuje $f(x)=|x|$.

\begin{wniosek}
\label{ma1:cor:inverse-lipschitz-mvt}
Jeśli przy tych samych założeniach $f'(x)\ge m>0$ w całym wnętrzu przedziału $I$, to $f$ jest ściśle rosnąca, a jej odwrotność na $f(I)$ jest funkcją Lipschitza ze stałą $1/m$.
\end{wniosek}
\begin{proof}
Dla $u<v$ stosujemy twierdzenie Lagrange'a, uzyskując $f(v)-f(u)=f'(c)(v-u)\ge m(v-u)>0$. W szczególności $f$ jest iniektywna. Dla dowolnych $s,t\in f(I)$ bierzemy $u=f^{-1}(s)$, $v=f^{-1}(t)$ i z powyższego oszacowania (ewentualnie po zamianie $u,v$) otrzymujemy $m|u-v|\le |s-t|$. Dzieląc przez $m>0$, dostajemy żądaną nierówność. Nie zakładamy istnienia górnego ograniczenia pochodnej.
\end{proof}
\begin{przyklad}
\label{ma1:ex:mvt-elementary-bounds}
Z oszacowań $|(\sin x)'|\le1$ i $|(\cos x)'|\le1$ wynika
$|\sin x-\sin y|\le|x-y|$ oraz $|\cos x-\cos y|\le|x-y|$ dla $x,y\in\R$.
Dla $0<a\le b$ i $x,y\in[a,b]$ pochodna logarytmu należy do $[1/b,1/a]$, zatem
\[
 \frac1b|x-y|\le|\log x-\log y|\le\frac1a|x-y|.
\]
Dla $x>0$ twierdzenie Lagrange'a zastosowane do $\log$ na $[1,1+x]$ daje $\log(1+x)=x/c$ dla pewnego $1<c<1+x$, więc $x/(1+x)<\log(1+x)<x$. Nierówności są ścisłe, ponieważ $c$ jest punktem wewnętrznym.
\end{przyklad}

\section{Przedłużenie pochodnej do końca przedziału}
\begin{twierdzenie}
\label{ma1:thm:derivative-endpoint-limit}
Niech $f:[a,b]\to\R$, $a<b$, będzie ciągła w $a$ i różniczkowalna na $(a,b)$. Jeżeli istnieje granica $\lim_{x\to a^+}f'(x)=L$ (skończona lub niewłaściwa), to
\[
 \lim_{h\to0^+}\frac{f(a+h)-f(a)}h=L.
\]
Analogiczne twierdzenie zachodzi na prawym końcu, przy ciągłości w $b$ i granicy pochodnej z lewej strony.
\end{twierdzenie}
\begin{proof}
Dla $0<h<b-a$ funkcja jest ciągła na $[a,a+h]$: w $a$ z założenia, w pozostałych punktach z różniczkowalności. Z twierdzenia Lagrange'a istnieje $c_h\in(a,a+h)$ takie, że iloraz różnicowy jest równy $f'(c_h)$. Mamy $0<c_h-a<h$, więc $c_h\to a^+$ przy $h\to0^+$. Z definicji granicy pochodnej otrzymujemy tezę, również gdy $L=+\infty$ lub $L=-\infty$. Dla prawego końca i $0<h<b-a$ stosujemy twierdzenie Lagrange'a do $[b-h,b]$: iloraz $(f(b)-f(b-h))/h$ jest równy $f'(d_h)$ z $b-h<d_h<b$. Stąd $d_h\to b^-$, a iloraz jest właśnie lewostronnym ilorazem różnicowym w $b$ dla przyrostu $-h$.
\end{proof}
Gdy $L$ jest skończona, otrzymujemy pochodną jednostronną równą $L$. Gdy granica jest niewłaściwa, iloraz różnicowy nie ma skończonej granicy. Odwrotność twierdzenia jest fałszywa: przykład \ref{ma1:ex:discontinuous-derivative} ma pochodną w $0$, ale jej wartości poza zerem nie mają nawet prawostronnej granicy w $0$.
